% ============================================================
%  论文框架结构图 v2 — 统一宽度分层架构
%  修复要点：
%    1. 所有层 fit box 统一 minimum width → 左右边界完全对齐
%    2. 所有模块相对于 x=0 对称排列 → box 中心永远在 x=0
%    3. 箭头从 box.south 到 box.north，全在 x=0 上 → 永远垂直
%    4. 方框高度由 fit 自动适应内容
%
%  编译命令：xelatex -interaction=nonstopmode example-paper-structure.tex
% ============================================================

\documentclass[12pt,a4paper]{ctexart}

\usepackage{geometry}
\geometry{left=1.5cm,right=1.5cm,top=2cm,bottom=2cm}

\usepackage{fontspec}
\usepackage{tikz}
\usetikzlibrary{positioning, fit, backgrounds, arrows.meta, calc}

% ============================================================
% 蓝色主题配色
% ============================================================
\definecolor{PaperBorder}{HTML}{B0C4DE}
\definecolor{PaperFill}{HTML}{FFFFFF}
\definecolor{PaperTitle}{HTML}{2C3E50}
\definecolor{PaperMuted}{HTML}{7F8C8D}
\definecolor{AccentOrange}{HTML}{D35400}
\definecolor{PaperLayerBg}{HTML}{F5F8FA}
\definecolor{PaperLayerBorder}{HTML}{D8E3EE}

% ============================================================
% 字体
% ============================================================
\setmainfont{Arial}
\setsansfont{Arial}
\setmonofont{Consolas}[Scale=0.90]
\setCJKsansfont{Microsoft YaHei}
\setCJKmainfont{Microsoft YaHei}

% ============================================================
% 文本样式命令
% ============================================================
\newcommand{\modtitle}{\small\sffamily\bfseries\color{PaperTitle}}
\newcommand{\moddesc}{\footnotesize\sffamily\color{PaperMuted}}

% ============================================================
% 统一宽度常量
% ============================================================
\def\TotalW{14.0cm}     % 所有层 fit box 的统一宽度
\def\LayerGap{1.2cm}    % 层间距

\begin{document}

\begin{center}
  {\Large\sffamily\bfseries\textcolor{PaperTitle}{论文框架结构图}}\\[0.3em]
  {\small\sffamily\itshape\textcolor{PaperMuted}{边缘计算环境下基于Flink的实时数据流处理与动态资源调度优化}}
\end{center}

\vspace{0.8em}

\begin{figure}[htbp]
\centering
\begin{tikzpicture}[
    >=Stealth,
    line width = 0.7pt,
    every node/.style = {inner sep=0pt, outer sep=0pt},
    % ---- 模块 ----
    module/.style = {
        rectangle,
        draw = PaperBorder,
        fill = PaperFill,
        text = PaperTitle,
        line width = 0.8pt,
        font = \small\sffamily,
        align = center,
        rounded corners = 5pt,
        inner sep = 8pt,
    },
    % ---- 层容器 ----
    layer-box/.style = {
        draw = PaperLayerBorder,
        fill = PaperLayerBg,
        line width = 0.7pt,
        dashed,
        dash pattern = on 4pt off 3pt,
        inner sep = 12pt,
        rounded corners = 8pt,
    },
    layer-title/.style = {
        font = \normalsize\sffamily\bfseries,
        text = PaperTitle,
        anchor = north west,
        inner sep = 0pt,
    },
    layer-tag/.style = {
        font = \footnotesize\sffamily\itshape,
        text = PaperMuted,
        anchor = north east,
        inner sep = 0pt,
    },
    % ---- 箭头 ----
    arr-down/.style = {
        ->, >=Stealth, thick, PaperBorder,
    },
    arr-label/.style = {
        font = \footnotesize\sffamily\bfseries,
        fill = white,
        text = PaperMuted,
        inner sep = 2pt,
        outer sep = 0pt,
    },
]

    % ============================================================
    %  层 0：论文概述（单模块，中心 x=0）
    % ============================================================
    \node[module, text width=12.8cm] (L0a) at (0, 0) {
        {\modtitle 论文标题}\\[2pt]
        {\moddesc 作者 / 摘要 / 关键词}
    };
    \begin{scope}[on background layer]
        \node[layer-box, fit=(L0a), minimum width=\TotalW] (box0) {};
    \end{scope}
    \node[layer-title] at ($(box0.north west)+(0,4pt)$) {论文概述};
    \node[layer-tag] at ($(box0.north east)+(0,4pt)$) {Title \& Abstract};

    % ============================================================
    %  层 1：引言（单模块）
    % ============================================================
    \coordinate (y1) at ($(box0.south)+(0,-\LayerGap)$);
    \node[module, text width=12.8cm, anchor=north] (L1a) at (y1) {
        {\modtitle 一、引言}\\[2pt]
        {\moddesc 研究背景 · 问题定义}\\[1pt]
        {\moddesc 研究目标 · 论文组织}
    };
    \begin{scope}[on background layer]
        \node[layer-box, fit=(L1a), minimum width=\TotalW] (box1) {};
    \end{scope}
    \node[layer-title] at ($(box1.north west)+(0,4pt)$) {第一章};
    \node[layer-tag] at ($(box1.north east)+(0,4pt)$) {Introduction};

    % ============================================================
    %  层 2：系统架构（双模块，中心 ±3.5）
    % ============================================================
    \coordinate (y2) at ($(box1.south)+(0,-\LayerGap)$);
    \node[module, text width=5.9cm, anchor=north] (L2a) at ($(y2)+(-3.5,0)$) {
        {\modtitle 2.1 云-边-端架构}\\[2pt]
        {\moddesc 三层协同模型}\\[1pt]
        {\moddesc 边缘节点职责}
    };
    \node[module, text width=5.9cm, anchor=north] (L2b) at ($(y2)+(3.5,0)$) {
        {\modtitle 2.2 Flink 引擎}\\[2pt]
        {\moddesc 流处理核心机制}\\[1pt]
        {\moddesc 窗口与状态管理}
    };
    \begin{scope}[on background layer]
        \node[layer-box, fit=(L2a)(L2b), minimum width=\TotalW] (box2) {};
    \end{scope}
    \node[layer-title] at ($(box2.north west)+(0,4pt)$) {第二章：系统架构};
    \node[layer-tag] at ($(box2.north east)+(0,4pt)$) {System Architecture};

    % ============================================================
    %  层 3：调度策略（双模块）
    % ============================================================
    \coordinate (y3) at ($(box2.south)+(0,-\LayerGap)$);
    \node[module, text width=5.9cm, anchor=north] (L3a) at ($(y3)+(-3.5,0)$) {
        {\modtitle 3.1 问题建模}\\[2pt]
        {\moddesc 多目标优化模型}\\[1pt]
        {\moddesc 约束条件分析}
    };
    \node[module, text width=5.9cm, anchor=north] (L3b) at ($(y3)+(3.5,0)$) {
        {\modtitle 3.2 算法设计}\\[2pt]
        {\moddesc 启发式调度策略}\\[1pt]
        {\moddesc 自适应参数调节}
    };
    \begin{scope}[on background layer]
        \node[layer-box, fit=(L3a)(L3b), minimum width=\TotalW] (box3) {};
    \end{scope}
    \node[layer-title] at ($(box3.north west)+(0,4pt)$) {第三章：调度策略};
    \node[layer-tag] at ($(box3.north east)+(0,4pt)$) {Scheduling Strategy};

    % ============================================================
    %  层 4：实验评估（四模块，中心 -5.25, -1.75, 1.75, 5.25）
    % ============================================================
    \coordinate (y4) at ($(box3.south)+(0,-\LayerGap)$);
    \node[module, text width=2.4cm, anchor=north] (L4a) at ($(y4)+(-5.25,0)$) {
        {\modtitle 4.1}\\[1pt]
        {\moddesc 实验环境}
    };
    \node[module, text width=2.4cm, anchor=north] (L4b) at ($(y4)+(-1.75,0)$) {
        {\modtitle 4.2}\\[1pt]
        {\moddesc 吞吐量对比}
    };
    \node[module, text width=2.4cm, anchor=north] (L4c) at ($(y4)+(1.75,0)$) {
        {\modtitle 4.3}\\[1pt]
        {\moddesc 延迟分析}
    };
    \node[module, text width=2.4cm, anchor=north] (L4d) at ($(y4)+(5.25,0)$) {
        {\modtitle 4.4}\\[1pt]
        {\moddesc 资源利用率}
    };
    \begin{scope}[on background layer]
        \node[layer-box, fit=(L4a)(L4b)(L4c)(L4d), minimum width=\TotalW] (box4) {};
    \end{scope}
    \node[layer-title] at ($(box4.north west)+(0,4pt)$) {第四章：实验评估};
    \node[layer-tag] at ($(box4.north east)+(0,4pt)$) {Experiments};

    % ============================================================
    %  层 5：总结（双模块）
    % ============================================================
    \coordinate (y5) at ($(box4.south)+(0,-\LayerGap)$);
    \node[module, text width=5.9cm, anchor=north] (L5a) at ($(y5)+(-3.5,0)$) {
        {\modtitle 五、相关工作}\\[2pt]
        {\moddesc 边缘计算研究现状}\\[1pt]
        {\moddesc 流处理调度方案}
    };
    \node[module, text width=5.9cm, anchor=north] (L5b) at ($(y5)+(3.5,0)$) {
        {\modtitle 六、结论与展望}\\[2pt]
        {\moddesc 研究总结}\\[1pt]
        {\moddesc 未来方向}
    };
    \begin{scope}[on background layer]
        \node[layer-box, fit=(L5a)(L5b), minimum width=\TotalW] (box5) {};
    \end{scope}
    \node[layer-title] at ($(box5.north west)+(0,4pt)$) {总结};
    \node[layer-tag] at ($(box5.north east)+(0,4pt)$) {Related Work \& Conclusion};

    % ============================================================
    %  层间箭头（全部在 x=0 垂直线上，永远不会歪）
    % ============================================================
    \draw[arr-down] (box0.south) -- (box1.north)
        node[arr-label, midway] {引入};
    \draw[arr-down] (box1.south) -- (box2.north)
        node[arr-label, midway] {架构设计};
    \draw[arr-down] (box2.south) -- (box3.north)
        node[arr-label, midway] {策略优化};
    \draw[arr-down] (box3.south) -- (box4.north)
        node[arr-label, midway] {验证};
    \draw[arr-down] (box4.south) -- (box5.north)
        node[arr-label, midway] {总结};

\end{tikzpicture}
\caption{\sffamily 论文框架结构图}
\label{fig:paper-structure}
\end{figure}

\end{document}
